% =============================================================================
%  cours_java_arrays_strings_lists.tex
%  Cours Java -- Arrays, Strings, List et Java Collections Framework
%  (interfaces et hierarchie des listes)
%
%  Compilation (deux passes pour la table des matieres) :
%     pdflatex cours_java_arrays_strings_lists.tex
%     pdflatex cours_java_arrays_strings_lists.tex
% =============================================================================
\documentclass[11pt,a4paper]{article}

\usepackage[utf8]{inputenc}
\usepackage[T1]{fontenc}
\usepackage[french]{babel}
\usepackage[margin=2.2cm]{geometry}
\usepackage[table]{xcolor}
\usepackage{listings}
\usepackage{tcolorbox}
\usepackage{enumitem}
\usepackage{titlesec}
\usepackage{longtable}
\usepackage{array}
\usepackage{booktabs}
\usepackage{fancyhdr}
\usepackage{hyperref}
\usepackage{amssymb}

% --- Couleurs (meme palette que le cours d'algorithmes) --------------------
\definecolor{primary}{HTML}{1F3A5F}
\definecolor{accent}{HTML}{2E86DE}
\definecolor{codebg}{HTML}{F4F7FB}
\definecolor{codekw}{HTML}{1F3A5F}
\definecolor{codestr}{HTML}{27AE60}
\definecolor{codecom}{HTML}{7F8C99}
\definecolor{rulecol}{HTML}{D5DEE8}

% --- Liens -------------------------------------------------------------
\hypersetup{colorlinks=true, linkcolor=accent, urlcolor=accent, pdftitle={Cours Java -- Arrays, Strings, List, Collections}}

% --- En-tete / pied de page ---------------------------------------------
\pagestyle{fancy}
\fancyhf{}
\fancyhead[L]{\small\color{primary}Cours Java -- Arrays, Strings, List, Collections}
\fancyhead[R]{\small\color{codecom}\thepage}
\renewcommand{\headrulewidth}{0.4pt}

% --- Titres de section ---------------------------------------------------
\titleformat{\section}{\Large\bfseries\color{primary}}{\thesection.}{0.6em}{}
\titleformat{\subsection}{\large\bfseries\color{accent}}{}{0em}{}

% --- Style du code Java ---------------------------------------------------
\lstdefinestyle{javastyle}{
  language=Java,
  backgroundcolor=\color{codebg},
  basicstyle=\ttfamily\small,
  keywordstyle=\color{codekw}\bfseries,
  stringstyle=\color{codestr},
  commentstyle=\color{codecom}\itshape,
  numbers=none,
  frame=single,
  rulecolor=\color{rulecol},
  breaklines=true,
  showstringspaces=false,
  columns=fullflexible,
  keepspaces=true,
  aboveskip=8pt, belowskip=8pt,
  xleftmargin=6pt, xrightmargin=6pt,
  morekeywords={String,Integer,List,ArrayList,LinkedList,Map,HashMap,Set,HashSet,Collection,Collections,Arrays,Optional,var},
  literate=%
    {é}{{\'e}}1 {è}{{\`e}}1 {ê}{{\^e}}1 {à}{{\`a}}1 {ô}{{\^o}}1
    {î}{{\^i}}1 {ç}{{\c c}}1 {É}{{\'E}}1 {È}{{\`E}}1 {ù}{{\`u}}1,
}
\lstset{style=javastyle}

% --- Boites "Quand l'utiliser" / info -------------------------------------
\newtcolorbox{infobox}[1]{colback=codebg,colframe=accent,title=#1,
  fonttitle=\bfseries\color{white},boxrule=0.5pt,arc=2pt,left=6pt,right=6pt,top=4pt,bottom=4pt}

\newcommand{\complexite}[2]{%
  \noindent\textbf{Complexite :}\quad Temps \texttt{#1} \quad|\quad Espace \texttt{#2}\par\smallskip}

\newcommand{\leetcode}[1]{\noindent\textbf{Problemes types :} #1\par\medskip}

% =============================================================================
\begin{document}

% --- Page de titre ----------------------------------------------------------
\begin{titlepage}
\centering
\vspace*{3cm}
{\color{primary}\Huge\bfseries Cours Java\par}
\vspace{0.6cm}
{\color{accent}\LARGE Arrays, Strings, List \& Java Collections Framework\par}
\vspace{1.2cm}
{\large Les structures de donnees de base en Java :\\
tableaux, chaines de caracteres, listes, et la hierarchie des interfaces\\
du Collections Framework.\par}
\vfill
{\color{codecom}\large Cours a imprimer -- support d'entrainement\par}
\vspace{2cm}
\end{titlepage}

\tableofcontents
\newpage

% =============================================================================
\section{Arrays (tableaux) en Java}
% =============================================================================

Un tableau (\texttt{array}) est la structure de donnees la plus basique de Java :
une zone memoire \textbf{contigue} de \textbf{taille fixe}, dont tous les
elements ont le \textbf{meme type}. C'est le bloc de construction sur lequel
sont batis \texttt{ArrayList}, \texttt{String} (en interne, un tableau de
\texttt{char}/\texttt{byte}), et une bonne partie du Collections Framework.

\textbf{L'idee.} Parce que la taille est fixee a la creation et que les
elements sont contigus en memoire, l'acces a l'indice \texttt{i} se fait par un
simple calcul d'adresse (\texttt{base + i * taille\_element}) : c'est pour cela
que l'acces par indice est en \texttt{O(1)}. En contrepartie, on ne peut ni
agrandir ni retrecir un tableau apres sa creation.

\begin{infobox}{Declaration et initialisation}
\end{infobox}

\begin{lstlisting}
int[] nums = new int[5];              // tableau de 5 entiers, initialises a 0
int[] nums2 = {1, 2, 3, 4, 5};        // initialisation directe
int[] nums3 = new int[]{1, 2, 3};     // forme explicite equivalente

int[][] matrix = new int[3][4];       // tableau 2D : 3 lignes, 4 colonnes
int[][] grid = {{1, 2}, {3, 4}};      // initialisation directe d'une matrice

System.out.println(nums.length);      // longueur : propriete, PAS une methode
\end{lstlisting}

\begin{infobox}{Piege classique}
\texttt{nums.length} est un \textbf{champ}, pas une methode : pas de parentheses.
A l'inverse, pour une \texttt{String}, c'est \texttt{str.length()} avec
parentheses (une methode). Confondre les deux est l'erreur de syntaxe la plus
frequente chez les debutants en Java venant d'autres langages.
\end{infobox}

\subsection{Parcours d'un tableau}

\begin{lstlisting}
int[] nums = {10, 20, 30, 40};

// Boucle indexee classique
for (int i = 0; i < nums.length; i++) {
    System.out.println(nums[i]);
}

// Boucle "for-each" (enhanced for) : plus lisible si on n'a pas besoin de l'indice
for (int n : nums) {
    System.out.println(n);
}
\end{lstlisting}

\textbf{Reconnaitre quand utiliser quelle boucle.} Utilise le \texttt{for-each}
des que tu n'as besoin que de la \emph{valeur}. Reviens a la boucle indexee des
que tu as besoin de l'\textbf{indice} lui-meme (modifier le tableau en place,
comparer des positions, parcourir a l'envers, ou avancer de plusieurs pas a la
fois).

\subsection{La classe utilitaire \texttt{Arrays}}

Le tableau lui-meme n'a presque aucune methode ; toutes les operations
utiles passent par la classe statique \texttt{java.util.Arrays}.

\begin{lstlisting}
import java.util.Arrays;

int[] nums = {5, 3, 8, 1};

Arrays.sort(nums);                          // tri en place, O(n log n)
System.out.println(Arrays.toString(nums));  // affichage lisible : [1, 3, 5, 8]

int[] copie = Arrays.copyOf(nums, 6);       // copie avec nouvelle taille (padding a 0)
int[] sousTab = Arrays.copyOfRange(nums, 1, 3); // sous-tableau [indice 1, 3[

boolean egaux = Arrays.equals(nums, copie); // comparaison element par element
Arrays.fill(nums, 0);                       // remplit tout le tableau avec 0

// binarySearch exige un tableau DEJA TRIE, sinon resultat indefini
int idx = Arrays.binarySearch(nums, 5);
\end{lstlisting}

\complexite{O(1) acces par indice ; O(n) recherche lineaire ; O(n log n) tri}{O(n)}

\begin{infobox}{Piege classique}
Un tableau d'objets (\texttt{int[][]}, \texttt{Object[]}) n'a \textbf{pas} de
\texttt{toString()} lisible par defaut : \texttt{System.out.println(nums)}
affiche une adresse memoire du type \texttt{[I@1b6d3586}. Il faut toujours
passer par \texttt{Arrays.toString(nums)} (tableau 1D) ou
\texttt{Arrays.deepToString(matrix)} (tableau 2D).
\end{infobox}

\subsection{Convertir entre tableau et liste}

C'est une source frequente de bugs subtils en entretien.

\begin{lstlisting}
Integer[] boxed = {1, 2, 3};
List<Integer> vue = Arrays.asList(boxed);   // VUE sur le tableau, taille fixe !
// vue.add(4);  // UnsupportedOperationException : on ne peut pas l'agrandir

List<Integer> vraieListe = new ArrayList<>(Arrays.asList(boxed)); // copie modifiable

// Attention : Arrays.asList ne fonctionne pas comme attendu sur un tableau primitif
int[] primitifs = {1, 2, 3};
List<int[]> mauvais = Arrays.asList(primitifs); // liste d'UN SEUL element (le tableau lui-meme) !
\end{lstlisting}

\begin{infobox}{Piege classique}
\texttt{Arrays.asList(tableauPrimitif)} sur un \texttt{int[]} ne cree pas une
\texttt{List<Integer>} de 3 elements, mais une \texttt{List<int[]>} d'un seul
element (a cause de l'autoboxing qui ne s'applique qu'aux types objets).
Il faut un tableau de \texttt{Integer[]} (type objet) pour que la conversion se
comporte comme attendu, ou utiliser \texttt{Arrays.stream(primitifs).boxed()}.
\end{infobox}

\newpage
% =============================================================================
\section{Strings (chaines de caracteres) en Java}
% =============================================================================

En Java, \texttt{String} est un \textbf{objet immuable} : une fois cree, son
contenu ne peut jamais changer. Toute operation qui "modifie" une chaine
(\texttt{+}, \texttt{substring}, \texttt{replace}...) cree en realite une
\textbf{nouvelle} chaine et laisse l'originale intacte.

\textbf{L'idee.} L'immuabilite permet a la JVM de reutiliser en toute securite
la meme instance pour deux litteraux identiques (le \emph{String Pool}), et
rend les \texttt{String} intrinsequement thread-safe. Le prix a payer : creer
une chaine dans une boucle avec des concatenations \texttt{+=} recree un
nouvel objet a chaque iteration, ce qui devient couteux en \texttt{O(n\textsuperscript{2})}
sur \texttt{n} concatenations.

\begin{infobox}{Quand l'utiliser}
Toute donnee textuelle immuable : cles de map, comparaisons, parsing.
Des que tu dois construire une chaine \textbf{progressivement} (boucle,
concatenations repetees), pense \texttt{StringBuilder} plutot que \texttt{String}.
\end{infobox}

\subsection{Methodes essentielles}

\begin{lstlisting}
String s = "Hello, World!";

s.length();                 // 13 -- methode, PAS un champ (contrairement a un tableau)
s.charAt(0);                 // 'H'
s.substring(7);              // "World!"      (a partir de l'indice 7)
s.substring(7, 12);          // "World"       (indice 7 inclus, 12 exclu)
s.indexOf("World");          // 7             (-1 si non trouve)
s.toLowerCase();              // "hello, world!"
s.toUpperCase();              // "HELLO, WORLD!"
s.trim();                     // supprime les espaces en debut/fin
s.strip();                    // equivalent moderne (Java 11+), gere aussi l'Unicode
s.replace("World", "Java");  // "Hello, Java!"
s.split(", ");                // String[] {"Hello", "World!"}
s.contains("World");          // true
s.startsWith("Hello");        // true
s.equals("hello, world!");    // false -- sensible a la casse
s.equalsIgnoreCase("HELLO, WORLD!"); // true

String.join("-", "a", "b", "c");  // "a-b-c"
"  ".isBlank();                    // true (Java 11+)
"".isEmpty();                       // true
\end{lstlisting}

\begin{infobox}{Piege classique}
Ne \textbf{jamais} comparer deux \texttt{String} avec \texttt{==} : cela
compare des references memoire, pas le contenu. \texttt{new String("abc") ==
"abc"} vaut \texttt{false}. Toujours utiliser \texttt{.equals()} (ou
\texttt{.equalsIgnoreCase()}).
\end{infobox}

\subsection{Construire une chaine efficacement : \texttt{StringBuilder}}

\begin{lstlisting}
StringBuilder sb = new StringBuilder();
for (int i = 0; i < 5; i++) {
    sb.append(i).append(",");     // append est chainable
}
sb.deleteCharAt(sb.length() - 1); // retire la derniere virgule
String resultat = sb.toString();  // "0,1,2,3,4"

sb.reverse();                       // inverse le contenu en place
sb.insert(0, "prefix-");            // insertion a un indice donne
\end{lstlisting}

\complexite{O(1) amorti par append}{O(n)}

\begin{infobox}{Reconnaitre le pattern}
Des que l'enonce implique de construire une chaine \emph{caractere par
caractere} ou \emph{morceau par morceau} dans une boucle (inverser une chaine,
construire un resultat progressif, palindromes), utilise \texttt{StringBuilder}
: cela transforme un cout quadratique en cout lineaire.
\end{infobox}

\subsection{Manipuler les caracteres}

\begin{lstlisting}
char c = 'a';
Character.isDigit(c);       // false
Character.isLetter(c);      // true
Character.isUpperCase(c);   // false
Character.toUpperCase(c);   // 'A'

// Convertir String en tableau de char pour un acces/tri caractere par caractere
char[] chars = "hello".toCharArray();
Arrays.sort(chars);                         // tri des caracteres
String trie = new String(chars);            // reconversion en String
\end{lstlisting}

\leetcode{Valid Anagram (LC 242), Reverse String (LC 344), Longest Common Prefix (LC 14), Valid Palindrome (LC 125).}

\newpage
% =============================================================================
\section{List en Java}
% =============================================================================

\texttt{List<E>} est une \textbf{interface} du Collections Framework
representant une sequence ordonnee d'elements, avec doublons autorises et
acces par indice. Elle resout la principale limite des tableaux : la
\textbf{taille fixe}.

\textbf{L'idee.} On programme presque toujours contre l'interface
\texttt{List<E>} plutot que contre une implementation concrete
(\texttt{ArrayList}, \texttt{LinkedList}) : cela permet de changer
d'implementation plus tard sans modifier le reste du code, et c'est la
convention Java standard (« coder contre une interface »).

\begin{lstlisting}
List<Integer> list = new ArrayList<>();   // le type declare est l'interface,
                                           // le type instancie est l'implementation

list.add(10);
list.add(20);
list.add(0, 5);           // insertion a l'indice 0 : [5, 10, 20]
list.get(1);              // 10
list.set(1, 99);          // remplace l'element a l'indice 1 : [5, 99, 20]
list.remove(0);            // remove(int) : supprime PAR INDICE -> [99, 20]
list.remove(Integer.valueOf(99)); // remove(Object) : supprime PAR VALEUR -> [20]
list.contains(20);         // true
list.indexOf(20);          // 0
list.size();                // 1
list.isEmpty();             // false
list.clear();               // vide la liste

List<Integer> literal = List.of(1, 2, 3); // liste IMMUABLE (Java 9+)
\end{lstlisting}

\begin{infobox}{Piege classique}
Sur une \texttt{List<Integer>}, \texttt{list.remove(0)} et
\texttt{list.remove(Integer.valueOf(0))} \textbf{ne font pas la meme chose} :
la premiere supprime l'element a l'\emph{indice} 0 (surcharge
\texttt{remove(int)}), la seconde supprime la \emph{valeur} 0 (surcharge
\texttt{remove(Object)}). C'est une ambiguite classique qui piege beaucoup de
candidats en entretien.
\end{infobox}

\subsection{\texttt{ArrayList} vs \texttt{LinkedList}}

Les deux implementent \texttt{List<E>}, mais avec des structures internes tres
differentes.

\begin{longtable}{p{3.2cm}p{5.4cm}p{5.4cm}}
\toprule
\textbf{Operation} & \textbf{ArrayList} (tableau redimensionnable) & \textbf{LinkedList} (liste doublement chainee) \\
\midrule
\endhead
Acces par indice \texttt{get(i)} & \texttt{O(1)} & \texttt{O(n)} \\
Ajout/suppression en fin & \texttt{O(1)} amorti & \texttt{O(1)} \\
Ajout/suppression en debut & \texttt{O(n)} (decalage) & \texttt{O(1)} \\
Ajout/suppression au milieu & \texttt{O(n)} (decalage) & \texttt{O(n)} pour trouver + \texttt{O(1)} pour inserer \\
Empreinte memoire & Compacte & Plus lourde (pointeurs prev/next par noeud) \\
\bottomrule
\end{longtable}

\begin{infobox}{Reconnaitre quand utiliser quoi}
\texttt{ArrayList} est le choix par defaut dans l'immense majorite des cas :
les acces par indice et les parcours dominent en pratique. Ne choisis
\texttt{LinkedList} que si tu fais \emph{beaucoup} d'insertions/suppressions
en tete ou en queue (elle implemente aussi \texttt{Deque<E>}, utile pour une
pile ou une file).
\end{infobox}

\begin{lstlisting}
LinkedList<Integer> deque = new LinkedList<>();
deque.addFirst(1);   // ajout en tete, O(1)
deque.addLast(2);     // ajout en queue, O(1)
deque.removeFirst();  // retrait en tete, O(1)
deque.peekLast();     // consultation sans retrait
\end{lstlisting}

\subsection{Parcourir et trier une liste}

\begin{lstlisting}
List<Integer> nums = new ArrayList<>(List.of(5, 3, 8, 1));

for (int n : nums) {                 // for-each
    System.out.println(n);
}

Collections.sort(nums);              // tri croissant en place, O(n log n)
nums.sort(Comparator.reverseOrder()); // tri decroissant, methode d'instance (Java 8+)

// Tri avec comparateur personnalise (ex : liste de String par longueur)
List<String> mots = new ArrayList<>(List.of("banane", "kiwi", "pomme"));
mots.sort((a, b) -> a.length() - b.length());
mots.sort(Comparator.comparingInt(String::length));  // equivalent, plus idiomatique
\end{lstlisting}

\begin{infobox}{Piege classique}
Modifier une liste (\texttt{add}/\texttt{remove}) \textbf{pendant} un
\texttt{for-each} leve une \texttt{ConcurrentModificationException}. Pour
supprimer des elements en parcourant, utilise un \texttt{Iterator} explicite
avec \texttt{iterator.remove()}, ou \texttt{list.removeIf(condition)}.
\end{infobox}

\begin{lstlisting}
// Suppression sure pendant un parcours
Iterator<Integer> it = nums.iterator();
while (it.hasNext()) {
    if (it.next() % 2 == 0) {
        it.remove();          // seule facon sure de supprimer pendant l'iteration
    }
}

// Equivalent moderne et plus lisible (Java 8+)
nums.removeIf(n -> n % 2 == 0);
\end{lstlisting}

\complexite{Voir tableau ci-dessus selon l'implementation}{O(n)}
\leetcode{Rotate Array (LC 189), Remove Duplicates from Sorted Array (LC 26), Merge Sorted Array (LC 88).}

\newpage
% =============================================================================
\section{Java Collections Framework : interfaces et hierarchie}
% =============================================================================

Le Collections Framework est l'ensemble unifie d'interfaces et de classes de
\texttt{java.util} pour manipuler des groupes d'objets. Comprendre sa
\textbf{hierarchie d'interfaces} est essentiel : en entretien comme en
production, on declare toujours des variables avec le type \textbf{interface}
le plus general possible (\texttt{List}, \texttt{Map}, \texttt{Set}), jamais
avec le type d'implementation concrete.

\subsection{Vue d'ensemble de la hierarchie}

\begin{lstlisting}
Iterable<E>
    |
Collection<E>  ------------------------------------
    |                    |                          |
  List<E>              Set<E>                    Queue<E>
    |                 /      \                       |
ArrayList        HashSet   TreeSet/            LinkedList / ArrayDeque
LinkedList       LinkedHashSet  SortedSet           |
                                                   Deque<E>

Map<K,V>  (NE derive PAS de Collection -- interface a part)
    |
  HashMap, LinkedHashMap, TreeMap (implemente SortedMap)
\end{lstlisting}

\begin{infobox}{Point souvent teste en entretien}
\texttt{Map<K,V>} ne fait \textbf{pas} partie de la hierarchie
\texttt{Collection<E>} : c'est une interface racine independante, car une map
manipule des \emph{paires} cle-valeur et non de simples elements. C'est une
question de comprehension frequemment posee pour verifier qu'on connait la
hierarchie plutot que de l'avoir simplement memorisee par cœur.
\end{infobox}

\subsection{L'interface \texttt{Collection<E>}}

Racine commune de \texttt{List}, \texttt{Set} et \texttt{Queue}. Elle definit
le contrat partage par toutes les collections d'elements uniques (au sens
"un seul type d'element", doublons ou non selon la sous-interface) :
\texttt{add}, \texttt{remove}, \texttt{contains}, \texttt{size},
\texttt{isEmpty}, \texttt{iterator}, \texttt{clear}, \texttt{stream}.

\subsection{\texttt{List<E>} : sequence ordonnee, doublons autorises}

Deja detaillee section precedente. Implementations : \texttt{ArrayList}
(tableau redimensionnable), \texttt{LinkedList} (liste chainee).

\subsection{\texttt{Set<E>} : elements uniques, pas de doublons}

\begin{longtable}{p{3.4cm}p{4.6cm}p{6cm}}
\toprule
\textbf{Implementation} & \textbf{Structure interne} & \textbf{Caracteristique} \\
\midrule
\endhead
\texttt{HashSet} & Table de hachage (adossee a une \texttt{HashMap}) & Pas d'ordre garanti ; \texttt{O(1)} moyen pour add/contains \\
\texttt{LinkedHashSet} & Table de hachage + liste chainee & Conserve l'ordre d'insertion ; \texttt{O(1)} moyen \\
\texttt{TreeSet} & Arbre rouge-noir & Elements toujours tries ; \texttt{O(log n)} pour add/contains \\
\bottomrule
\end{longtable}

\begin{lstlisting}
Set<Integer> set = new HashSet<>();
set.add(5);
set.add(5);              // ignore silencieusement, deja present
set.contains(5);          // true, O(1) moyen

Set<Integer> ordonne = new TreeSet<>(Set.of(5, 1, 3)); // {1, 3, 5} -- toujours trie
\end{lstlisting}

\begin{infobox}{Reconnaitre le pattern}
Des que l'enonce parle d'\emph{unicite}, de \emph{doublons a supprimer}, ou
de tester une appartenance repetee dans une boucle (\texttt{if (list.contains(x))}
en \texttt{O(n)} a chaque appel), remplace la structure par un \texttt{Set} :
le test d'appartenance passe de \texttt{O(n)} a \texttt{O(1)} moyen.
\end{infobox}

\subsection{\texttt{Queue<E>} et \texttt{Deque<E>}}

\texttt{Queue<E>} modelise une file (FIFO -- premier entre, premier sorti).
\texttt{Deque<E>} (\emph{double-ended queue}) generalise a une file a deux
extremites, et peut donc aussi servir de \textbf{pile} (LIFO).

\begin{lstlisting}
Queue<Integer> file = new LinkedList<>();
file.offer(1);           // ajoute en fin (equivalent add, mais ne leve pas d'exception si echec)
file.poll();               // retire et retourne la tete (null si vide, contrairement a remove())
file.peek();                // consulte la tete sans retirer

Deque<Integer> pile = new ArrayDeque<>();  // ArrayDeque : implementation recommandee pour une pile
pile.push(1);              // empile
pile.push(2);
pile.pop();                 // depile -> 2 (LIFO)
pile.peek();                 // consulte le sommet -> 1
\end{lstlisting}

\begin{infobox}{Piege classique}
La classe historique \texttt{java.util.Stack} existe toujours mais est
\textbf{deconseillee} : elle herite de \texttt{Vector}, une classe
synchronisee (donc plus lente inutilement) et mal concue. Utiliser
\texttt{ArrayDeque} pour implementer une pile en Java moderne.
\end{infobox}

\subsection{Comparaison synthetique des trois familles}

\begin{longtable}{p{2.6cm}p{3.4cm}p{4.2cm}p{4.2cm}}
\toprule
\textbf{Interface} & \textbf{Doublons ?} & \textbf{Ordre ?} & \textbf{Cas d'usage typique} \\
\midrule
\endhead
\texttt{List<E>} & Oui & Ordre d'insertion, indexe & Sequence, historique, resultat ordonne \\
\texttt{Set<E>} & Non & Depend de l'implementation & Unicite, appartenance rapide \\
\texttt{Queue}/\texttt{Deque} & Oui & FIFO ou LIFO & File d'attente, pile, BFS/DFS \\
\texttt{Map<K,V>} & Cles non, valeurs oui & Depend de l'implementation & Association cle $\to$ valeur, comptage, cache \\
\bottomrule
\end{longtable}

\begin{infobox}{A retenir}
Toujours \textbf{declarer le type interface} et \textbf{instancier
l'implementation} : \texttt{List<Integer> l = new ArrayList<>();} plutot que
\texttt{ArrayList<Integer> l = new ArrayList<>();}. Cela decouple le code du
choix d'implementation et permet de le changer plus tard (par ex. passer
d'\texttt{ArrayList} a \texttt{LinkedList}) sans toucher au reste du programme.
\end{infobox}

\leetcode{Design HashSet (LC 705), Implement Queue using Stacks (LC 232), Min Stack (LC 155).}

\end{document}
